% GENERATED FILE. Edit canonical lessons, then run npm run generate:books.
% canonical-source-hash: fnv1a64:24347318fdd730aa
% canonical-lessons: MR-C188-dindarsika, MR-C188-khekda, MR-C188-garud, MR-C188-kokila, MR-C188-madhmasi, MR-R188-first-pass-words-from-chapters-180-183, MR-R188-second-pass-words-from-chapters-184-188

\chapter{Calendar, Crab, Eagle, Cuckoo, Bee}
\label{ch:mr-calendar-crab-eagle-cuckoo-bee}
\hlchaptermodality{188}

\begin{chapteropening}
\textbf{By the end of this chapter:} \emph{I can say a calendar, a crab, an eagle, a cuckoo and a bee.}

Complete the last lesson of chapter 188: I can say a calendar, a crab, an eagle, a cuckoo and a bee.
\end{chapteropening}

\section[\texorpdfstring{\mr{दिनदर्शिका} (dindarśikā)}{दिनदर्शिका (dindarśikā)}]{\mr{दिनदर्शिका} (dindarśikā) — a calendar}

\label{lesson:MR-C188-dindarsika}



\begin{quote}
Before the new one: say the Marathi for a piece of jewellery, then the Marathi for a bangle.
\end{quote}

\subsection*{You'll want to know: \mr{दिनदर्शिका}}
\textbf{\mr{दिनदर्शिका}} — \emph{dindarśikā} — \textquotedblleft{}a calendar\textquotedblright{}.

\begin{grammarlens}[title={What you've built}]
One more word for this chapter.
\end{grammarlens}

\subsection*{Your turn}
\begin{itemize}
\raggedright
  \item \emph{Say it:} \emph{dindarśikā}
  \item \emph{Say it:} \emph{dindarśikā}, once more
  \item \emph{Say it:} say \emph{dindarśikā}
  \item \emph{Recall:} say the Marathi for a piece of jewellery, then the Marathi for a bangle, then say \emph{dindarśikā} again
\end{itemize}

\begin{tcolorbox}[breakable,colback=teal!4,colframe=teal!35!black,title={Before you move on}]
What does \mr{दिनदर्शिका} mean? (\textquotedblleft{}A calendar\textquotedblright{}.) Say it once more.
\end{tcolorbox}
\section[\texorpdfstring{\mr{खेकडा} (khekḍā)}{खेकडा (khekḍā)}]{\mr{खेकडा} (khekḍā) — a crab}

\label{lesson:MR-C188-khekda}



\begin{quote}
Before the new one: say the Marathi for a bangle, then the Marathi for a calendar.
\end{quote}

\subsection*{You'll want to know: \mr{खेकडा}}
\textbf{\mr{खेकडा}} — \emph{khekḍā} — \textquotedblleft{}a crab\textquotedblright{}.

\begin{grammarlens}[title={What you've built}]
One more word for this chapter.
\end{grammarlens}

\subsection*{Your turn}
\begin{itemize}
\raggedright
  \item \emph{Say it:} \emph{khekḍā}
  \item \emph{Say it:} \emph{khekḍā}, once more
  \item \emph{Say it:} say \emph{khekḍā}
  \item \emph{Recall:} say the Marathi for a bangle, then the Marathi for a calendar, then say \emph{khekḍā} again
\end{itemize}

\begin{tcolorbox}[breakable,colback=teal!4,colframe=teal!35!black,title={Before you move on}]
What does \mr{खेकडा} mean? (\textquotedblleft{}A crab\textquotedblright{}.) Say it once more.
\end{tcolorbox}
\section[\texorpdfstring{\mr{गरुड} (garuḍ)}{गरुड (garuḍ)}]{\mr{गरुड} (garuḍ) — an eagle}

\label{lesson:MR-C188-garud}



\begin{quote}
Before the new one: say the Marathi for a calendar, then the Marathi for a crab.
\end{quote}

\subsection*{You'll want to know: \mr{गरुड}}
\textbf{\mr{गरुड}} — \emph{garuḍ} — \textquotedblleft{}an eagle\textquotedblright{}.

\begin{grammarlens}[title={What you've built}]
One more word for this chapter.
\end{grammarlens}

\subsection*{Your turn}
\begin{itemize}
\raggedright
  \item \emph{Say it:} \emph{garuḍ}
  \item \emph{Say it:} \emph{garuḍ}, once more
  \item \emph{Say it:} say \emph{garuḍ}
  \item \emph{Recall:} say the Marathi for a calendar, then the Marathi for a crab, then say \emph{garuḍ} again
\end{itemize}

\begin{tcolorbox}[breakable,colback=teal!4,colframe=teal!35!black,title={Before you move on}]
What does \mr{गरुड} mean? (\textquotedblleft{}An eagle\textquotedblright{}.) Say it once more.
\end{tcolorbox}
\section[\texorpdfstring{\mr{कोकिळा} (kokiḷā)}{कोकिळा (kokiḷā)}]{\mr{कोकिळा} (kokiḷā) — a cuckoo}

\label{lesson:MR-C188-kokila}



\begin{quote}
Before the new one: say the Marathi for a crab, then the Marathi for an eagle.
\end{quote}

\subsection*{You'll want to know: \mr{कोकिळा}}
\textbf{\mr{कोकिळा}} — \emph{kokiḷā} — \textquotedblleft{}a cuckoo\textquotedblright{}.

\begin{grammarlens}[title={What you've built}]
One more word for this chapter.
\end{grammarlens}

\subsection*{Your turn}
\begin{itemize}
\raggedright
  \item \emph{Say it:} \emph{kokiḷā}
  \item \emph{Say it:} \emph{kokiḷā}, once more
  \item \emph{Say it:} say \emph{kokiḷā}
  \item \emph{Recall:} say the Marathi for a crab, then the Marathi for an eagle, then say \emph{kokiḷā} again
\end{itemize}

\begin{tcolorbox}[breakable,colback=teal!4,colframe=teal!35!black,title={Before you move on}]
What does \mr{कोकिळा} mean? (\textquotedblleft{}A cuckoo\textquotedblright{}.) Say it once more.
\end{tcolorbox}
\section[\texorpdfstring{\mr{मधमाशी} (madhmāśī)}{मधमाशी (madhmāśī)}]{\mr{मधमाशी} (madhmāśī) — a bee}

\label{lesson:MR-C188-madhmasi}



\begin{quote}
Before the new one: say the Marathi for an eagle, then the Marathi for a cuckoo.
\end{quote}

\subsection*{You'll want to know: \mr{मधमाशी}}
\textbf{\mr{मधमाशी}} — \emph{madhmāśī} — \textquotedblleft{}a bee\textquotedblright{}.

\begin{grammarlens}[title={What you've built}]
One more word for this chapter.
\end{grammarlens}

\subsection*{Your turn}
\begin{itemize}
\raggedright
  \item \emph{Say it:} \emph{madhmāśī}
  \item \emph{Say it:} \emph{madhmāśī}, once more
  \item \emph{Say it:} say \emph{madhmāśī}
  \item \emph{Recall:} say the Marathi for an eagle, then the Marathi for a cuckoo, then say \emph{madhmāśī} again
\end{itemize}

\begin{tcolorbox}[breakable,colback=teal!4,colframe=teal!35!black,title={Before you move on}]
What does \mr{मधमाशी} mean? (\textquotedblleft{}A bee\textquotedblright{}.) Say it once more.
\end{tcolorbox}
\section[(dialogue)]{Words from chapters 180-183 — a first pass}

\label{lesson:MR-R188-first-pass-words-from-chapters-180-183}



\begin{quote}
Say the words for a cuckoo and a bee.
\end{quote}

\subsection*{Your turn}
Say these aloud:

\begin{itemize}
\raggedright
  \item a writer, a painter, a photographer, a journalist, an architect
  \item a daughter-in-law, a bridegroom, a bride, a widow, success
  \item a discount, a coin, a banknote, a loan, a tax
  \item a tunnel, a fort, a tower, a monument, a statue
\end{itemize}

\begin{tcolorbox}[breakable,colback=teal!4,colframe=teal!35!black,title={Before you move on}]
A writer? (\textbf{\mr{लेखक} / \mr{लेखिका}}, \emph{lekhak / lekhikā}.) A painter? (\textbf{\mr{चित्रकार}}, \emph{citrakār}.) A photographer? (\textbf{\mr{छायाचित्रकार}}, \emph{chāyācitrakār}.) A journalist? (\textbf{\mr{पत्रकार}}, \emph{patrakār}.) An architect? (\textbf{\mr{वास्तुविशारद}}, \emph{vāstuviśārad}.) A daughter-in-law? (\textbf{\mr{सून}}, \emph{sūn}.) A bridegroom? (\textbf{\mr{नवरदेव}}, \emph{navardev}.) A bride? (\textbf{\mr{नवरी}}, \emph{navrī}.) A widow? (\textbf{\mr{विधवा}}, \emph{vidhvā}.) Success? (\textbf{\mr{यश}}, \emph{yaś}.) A discount? (\textbf{\mr{सवलत}}, \emph{savlat}.) A coin? (\textbf{\mr{नाणे}}, \emph{nāṇe}.) A banknote? (\textbf{\mr{नोट}}, \emph{noṭ}.) A loan? (\textbf{\mr{कर्ज}}, \emph{karja}.) A tax? (\textbf{\mr{कर}}, \emph{kar}.) A tunnel? (\textbf{\mr{बोगदा}}, \emph{bogdā}.) A fort? (\textbf{\mr{किल्ला}}, \emph{killā}.) A tower? (\textbf{\mr{मनोरा}}, \emph{manorā}.) A monument? (\textbf{\mr{स्मारक}}, \emph{smārak}.) A statue? (\textbf{\mr{पुतळा}}, \emph{putaḷā}.)
\end{tcolorbox}
\section[(dialogue)]{Words from chapters 184-188 — a second pass}

\label{lesson:MR-R188-second-pass-words-from-chapters-184-188}



\begin{quote}
Say the words for smoke and a flame.
\end{quote}

\subsection*{Your turn}
Say these aloud:

\begin{itemize}
\raggedright
  \item smoke, a flame, the climate, the environment, pollution
  \item a gym, cricket, a victory, a defeat, a competition
  \item a papad, a chutney, a salad, a semolina pudding, a modak
  \item trousers, socks, a handkerchief, a piece of jewellery, a bangle
  \item a calendar, a crab, an eagle, a cuckoo, a bee
\end{itemize}

\begin{tcolorbox}[breakable,colback=teal!4,colframe=teal!35!black,title={Before you move on}]
Smoke? (\textbf{\mr{धूर}}, \emph{dhūr}.) A flame? (\textbf{\mr{ज्योत}}, \emph{jyot}.) The climate? (\textbf{\mr{हवामान}}, \emph{havāmān}.) The environment? (\textbf{\mr{पर्यावरण}}, \emph{paryāvaraṇ}.) Pollution? (\textbf{\mr{प्रदूषण}}, \emph{pradūṣaṇ}.) A gym? (\textbf{\mr{व्यायामशाळा}}, \emph{vyāyāmśāḷā}.) Cricket? (\textbf{\mr{क्रिकेट}}, \emph{krikeṭ}.) A victory? (\textbf{\mr{विजय}}, \emph{vijay}.) A defeat? (\textbf{\mr{पराभव}}, \emph{parābhav}.) A competition? (\textbf{\mr{स्पर्धा}}, \emph{spardhā}.) A papad? (\textbf{\mr{पापड}}, \emph{pāpaḍ}.) A chutney? (\textbf{\mr{चटणी}}, \emph{caṭṇī}.) A salad? (\textbf{\mr{कोशिंबीर}}, \emph{kośimbīr}.) A semolina pudding? (\textbf{\mr{शिरा}}, \emph{śirā}.) A modak? (\textbf{\mr{मोदक}}, \emph{modak}.) Trousers? (\textbf{\mr{विजार}}, \emph{vijār}.) Socks? (\textbf{\mr{मोजे}}, \emph{moje}.) A handkerchief? (\textbf{\mr{रुमाल}}, \emph{rumāl}.) A piece of jewellery? (\textbf{\mr{दागिना}}, \emph{dāginā}.) A bangle? (\textbf{\mr{बांगडी}}, \emph{bāṅgḍī}.) A calendar? (\textbf{\mr{दिनदर्शिका}}, \emph{dindarśikā}.) A crab? (\textbf{\mr{खेकडा}}, \emph{khekḍā}.) An eagle? (\textbf{\mr{गरुड}}, \emph{garuḍ}.) A cuckoo? (\textbf{\mr{कोकिळा}}, \emph{kokiḷā}.) A bee? (\textbf{\mr{मधमाशी}}, \emph{madhmāśī}.)
\end{tcolorbox}
